\section{Related Work}

\paragraph{User-centric environments for LLM evaluation.}  
Recent work has increasingly emphasized the importance of evaluating LLMs under realistic, user-centric conditions. Traditional benchmarks assess task success assuming fully specified prompts, overlooking how users often communicate underspecified, incremental, or indirect goals. To address this, benchmarks like User Reported Scenarios \citep{wang2024user}, Intention-in-Interaction \citep{qian2024tell}, and WildBench \citep{lin2024wildbench} compile real user queries to evaluate whether models align with nuanced human intent. These datasets emphasize preference satisfaction, multi-intent understanding, and performance in ``in-the-wild'' conditions. Complementing this, other benchmarks such as MINT \citep{wang2024mint}, PrefEval \citep{zhao2025prefeval}, $\tau$-Bench \citep{yao2024tau}, and $\tau^{2}$-Bench~\citep{barres2025tau} focus on dynamic, multi-turn interactions, testing whether agents can incorporate feedback, handle evolving preferences, and maintain user alignment over time. However, several limitations remain. For example, a large portion of interactions are automatically synthesized in these benches, often resulting in overly lengthy and unnatural goal formulations that diverge from how users typically express themselves. In contrast, {\gymname} explicitly models the three core interactive traits of user communication that are critical for real-world alignment. It also spans hundreds of diverse goal configurations, enabling fine-grained evaluation on a wide range of user intents, thus providing a more comprehensive and behaviorally grounded setting for assessing user-centric agents. We provide a comprehensive comparison of our gym environment with recent environments and benchmarks in \Cref{tab:comparison}.

\begin{table*}[!t]
\centering
\resizebox{\linewidth}{!}{
\begin{tabular}{lccccccccc}
\toprule
\textbf{Benchmark} 
& \textbf{\makecell{Multi-turn\\Interaction}} 
& \textbf{\makecell{Goal\\Ambiguity}} 
& \textbf{\makecell{Tool\\Use}} 
& \textbf{\makecell{Dynamic\\State}} 
& \textbf{\makecell{Multi-Aspect\\Reasoning}} 
& \textbf{\makecell{User\\Simulation}} 
& \textbf{\makecell{Domain\\Diversity}} 
& \textbf{\makecell{Custom-\\izable}} 
& \textbf{\makecell{Scal-\\able}} \\
\midrule
\textit{$\tau^2$-Bench\citep{barres2025tau}}          
& \cmark  % dual-control dialogue
& \xmark  % goals are defined
& \cmark  % tool calls in telecom
& \cmark  % world state updates
& \xmark  % single domain (telecom)
& \cmark  % user acts in environment
& \xmark  % one domain
& \gmark  % partially configurable
& \gmark  % task generation is scalable
\\
\textit{$\tau$-Bench\citep{yao2024tau}}                
& \cmark  % multi-turn with simulated user
& \xmark  % goal is clearly specified
& \cmark  % tool usage required
& \cmark  % DB state changes
& \gmark  % multiple goals per domain
& \cmark  % simulated user
& \gmark  % across two real-world domains
& \gmark  % domain-specific but modular
& \gmark  % simulation allows scaling
\\
\textit{ToolSandbox\citep{lu2024toolsandbox}}              
& \cmark  % multi-turn with tool use
& \gmark  % includes insufficient info cases
& \cmark  % simulated tools
& \cmark  % state transitions via tools
& \gmark  % multi-tool, but still atomic tasks
& \cmark  % scripted simulated user
& \gmark  % across several plugins/tools
& \gmark  % API + plug-in design is modular
& \gmark  % framework allows scaling
\\
\textit{MINT\citep{wang2024mint}}                           
& \cmark  % iterative feedback loop
& \gmark  % open-ended tasks
& \cmark  % Python exec tool
& \cmark  % partial feedback changes state
& \gmark  % code, math, etc.
& \cmark  % simulated feedback agent
& \gmark  % math, code, reasoning
& \gmark  % tunable reward & feedback
& \gmark  % task remixing enabled
\\
\textit{IN3\citep{qian2024tell}}  
& \cmark  % multi-turn clarifications
& \cmark  % vague goals, implicit intent
& \xmark  % no tool use
& \xmark  % no changing state
& \xmark  % single goal types
& \cmark  % simulated users
& \gmark  % some variation in scenarios
& \xmark  % not modular/extensible
& \xmark  % curated, not scalable
\\
\textit{ToolTalk\citep{farn2023tooltalk}}                   
& \gmark  % mostly one-shot, but tools require inner planning
& \xmark  % user intent is clear
& \cmark  % broad tool use
& \xmark  % no evolving state
& \xmark  % usually one goal per task
& \xmark  % static user query
& \gmark  % multiple domains/tools
& \xmark  % no exposed framework
& \xmark  % small curated benchmark
\\
\textit{API-Bank\citep{li2023apibank}}                      
& \xmark  % single-turn prompts
& \xmark  % explicit goals
& \cmark  % tool (API) calling required
& \xmark  % no persistent state
& \xmark  % single-step planning
& \xmark  % no interaction
& \gmark  % many APIs/tasks
& \gmark  % released training + eval sets
& \gmark  % scalable data generation
\\
\textit{MetaTool\citep{huang2023metatool}}                  
& \xmark  % single prompt
& \xmark  % clear instructions
& \cmark  % tool use selection tested
& \xmark  % no evolving state
& \xmark  % no subtask coordination
& \xmark  % no user responses
& \gmark  % many tools/domains
& \gmark  % clean task schema
& \gmark  % 21k+ prompts generated
\\
\midrule
\textbf{{\gymname} (Ours)}              
& \cmark  % dialogue with vague user
& \cmark  % implicit, evolving intent
& \cmark  % tool-backed search
& \cmark  % stateful interaction
& \cmark  % multi-aspect: flight + hotel + etc.
& \cmark  % simulated user w/ heuristics
& \cmark  % various goals/preferences
& \cmark  % modular Gym API
& \cmark  % auto data generation + tiers
\\
\bottomrule
\end{tabular}
}
\caption{For each existing benchmark, the table indicates whether the corresponding trait is fully addressed (\cmark), partially addressed (\gmark), or not addressed (\xmark). We provide detailed explanation for each trait being compared in \Cref{sec:apdx_comparison}.}
\label{tab:comparison}
\end{table*}

\paragraph{User-centric agent designs.}  
Designing agents that collaborate effectively with users requires modeling ambiguity, evolving intent, and user-specific preferences. Standard instruction-tuned models often hallucinate intent or avoid asking clarifying questions. Recent work aims to teach agents to proactively clarify underspecified instructions. For instance, models trained with simulated ambiguity resolution \citep{zhang2024clarify, chen2025act} better recognize when to ask versus answer. Other work tackles user modeling directly: TravelPlanner+ \citep{singh2024personal} and PRELUDE \citep{gao2024prelude} build agents that personalize responses using either explicit profiles or latent preferences inferred from user edits. While these designs make progress toward adaptive, user-aware agents, they often target narrow personalization tasks or rely on static user models; in contrast, our work provides a general-purpose evaluation environment that systematically tests an agent’s ability to dynamically uncover and respond to emergent user intent across varied and customizable interaction scenarios.
